\documentclass[9pt,landscape,a4paper]{extarticle}
\usepackage[utf8]{inputenc}
\usepackage[vietnamese]{babel}
\usepackage[T1]{fontenc}
\usepackage{fontspec}
\setmainfont{Libertinus Serif} % As per user's font

% Layout and multicolumn
\usepackage[margin=0.8cm, top=1.5cm, bottom=1.5cm]{geometry}
\usepackage{multicol}
\setlength{\columnsep}{0.8cm}
\setlength{\columnseprule}{0.1pt}

% Math and symbols
\usepackage{amsmath,amssymb}
\usepackage{wasysym}
\usepackage{xcolor}

% Custom Colors (Stanford Cheatsheet style)
\definecolor{stanfordblue}{RGB}{1, 104, 149}
\definecolor{stanfordred}{RGB}{140, 21, 21}
\definecolor{stanfordgreen}{RGB}{0, 100, 0}
\definecolor{cyanbox}{RGB}{220, 245, 255}

% Section styling (mimicking the sheet)
\usepackage{titlesec}
\titleformat{\section}
  {\normalfont\normalsize\bfseries}{\thesection}{1em}{}
\titlespacing*{\section}{0pt}{1.5ex plus 1ex minus .2ex}{0.5ex plus .2ex}

\titleformat{\subsection}
  {\normalfont\small\bfseries}{\thesubsection}{1em}{}
\titlespacing*{\subsection}{0pt}{1ex plus 1ex minus .2ex}{0.5ex plus .2ex}

% Fancy header
\usepackage{fancyhdr}
\pagestyle{fancy}
\fancyhf{}
\renewcommand{\headrulewidth}{0.4pt}
\renewcommand{\footrulewidth}{0.4pt}

% Header content
\lhead{\textsc{MAT3562E} -- \textsc{COMPUTER VISION}}
\rhead{\colorbox{cyan!20}{\makebox[4cm][c]{\textbf{Bùi Quang Chiến -- 23001837}}}}

% Footer content
\lfoot{\textsc{HUS -- VNU}}
\cfoot{\thepage}
\rfoot{\textsc{SPRING 2026}}

% Dense lists
\usepackage{enumitem}
\setlist{nosep, leftmargin=1.2em}
\setlist[itemize,1]{label={\textcolor{stanfordgreen}{\footnotesize$\square$}}}
\setlist[itemize,2]{label=--}

\begin{document}
\begin{center}
    {\Large \bfseries VIP Cheatsheet: LAB08} \\
    \vspace{0.3cm}
    {\large \textsc{Bùi Quang Chiến}} \\
    \vspace{0.2cm}
    {\small April 2026}
\end{center}
\vspace{0.3cm}

\begin{multicols*}{3}

\begin{center}
    {\Large \bfseries VIP Cheatsheet: LAB08} \\
    \vspace{0.3cm}
    {\large \textsc{Bùi Quang Chiến}} \\
    \vspace{0.2cm}
    {\small April 2026}
\end{center}
\vspace{0.3cm}

\begin{multicols*}{3}

\begin{center}
  {\large\bfseries\color{hdr} LAB 08 CHEATSHEET --- Phân Đoạn \& Màu Cơ Bản} \\[1pt]
  {\small Bùi Quang Chiến \textbullet{} 23001837 \textbullet{} MAT3562E Computer Vision 2025--2026}
\end{center}
\hrule\vspace{2pt}

\begin{multicols}{3}

\section{1. Các Threshold trên Ảnh Màu}
\begin{itemize}
    \item \textbf{Otsu trên G/B/L* hoặc V}: Phù hợp khi vật thể sáng/tối phân minh.
    \item \textbf{Otsu trên S (Saturation)}: Rất hiệu quả nếu vật thể có màu nổi bật (colorful) trên phông nền xám/đơn sắc.
\end{itemize}
\begin{tipbox}
\textbf{Mô hình HSV 2 bước:}\\
1. Threshold trên S loại bỏ background xám.\\
2. Threshold trên V bên trong mask S để phân tầng độ sáng.
\end{tipbox}
\begin{tipbox}
\textbf{Lab Distance Threshold:}\\
$d(x,y) = \|c_{\mathrm{Lab}}(x,y) - c_{0,\mathrm{Lab}}\|_2$\\
Nghịch đảo Otsu hoặc fixed threshold $T$ trên ma trận khoảng cách $d$.
\end{tipbox}

\section{2. Dò Biên Màu (Color Edge Detection)}
\begin{itemize}
    \item \textbf{Đa kênh Canny}: Tách R,G,B $\to$ Canny từng kênh $\to$ Hợp (OR). Gây trùng viền và nhiều nhiễu mép.
    \item \textbf{Di Zenzo Gradient}: Mô hình Tensor.\\
    $g_{xx} = I_x^\top I_x$, $\quad g_{yy} = I_y^\top I_y$, $\quad g_{xy} = I_x^\top I_y$.\\
    Biên độ: $M = \sqrt{\frac{1}{2}\left(g_{xx} + g_{yy} + \sqrt{(g_{xx}-g_{yy})^2 + 4g_{xy}^2}\right)}$.\\
    Phát hiện được nét khác màu nhưng cùng chói!
    \item \textbf{Color Canny}: Dùng $M$ và $\theta$ Di Zenzo làm đầu vào cho NMS và Hysteresis $\to$ Cực xịn.
\end{itemize}

\section{3. Phân Đoạn (Segmentation)}
\begin{infobox}
\textbf{Region Growing (Phát triển vùng):}\\
Hoạt động cực mịn ở không gian Lab do sự tách biệt kênh chói. Metric L2 trong Lab tuân thủ nghiêm ngặt cảm nhận mắt người.
\end{infobox}
\begin{infobox}
\textbf{K-means Tọa độ (Spatial Coherence):}\\
Feature: $f(x,y) = [L^*, a^*, b^*, \lambda x, \lambda y]$. Trọng số $\lambda$ ép siêu điểm gần gũi đứng cạnh nhau phải về cùng 1 Cluster.
\end{infobox}
\begin{infobox}
\textbf{Split \& Merge:}\\
Chia \textbf{tứ phân} nếu Phương sai $\max_c \text{Var}_c(R) > T_{\sigma^2}$. Ghép lại lân cận nếu Mean tương đồng.
\end{infobox}

\end{multicols}

\end{multicols*}

\end{multicols*}
\end{document}
